% ---------------------------------------------------------------------------
% Proyecto I - Robot Aspiradora Autonomo (Yocto)
% CE-1113 Sistemas Empotrados - TEC
%
% Documento de Aprendizaje Continuo (AC): indicadores AC1 a AC4 del enunciado.
%
% Compilar:  pdflatex aprendizaje-continuo.tex   (dos veces)
% ---------------------------------------------------------------------------
\documentclass[11pt,a4paper]{article}

\usepackage[T1]{fontenc}
\usepackage[utf8]{inputenc}
\usepackage{lmodern}
\usepackage[spanish,es-noshorthands,es-noquoting]{babel}
\usepackage[margin=2.4cm]{geometry}
\usepackage{xcolor}
\usepackage{booktabs}
\usepackage{tabularx}
\usepackage{longtable}
\usepackage{array}
\usepackage{enumitem}
\usepackage[hidelinks]{hyperref}

\setlength{\parskip}{0.45em}
\setlength{\parindent}{0pt}
\emergencystretch=4em
\setlist{itemsep=1pt,topsep=2pt}
\renewcommand{\arraystretch}{1.18}

\definecolor{cFondo}{RGB}{245,247,250}
\definecolor{cPend}{RGB}{176,32,32}

% Identificadores y rutas con guion bajo, sin escaparlos a mano.
\newcommand{\id}[1]{\texttt{\detokenize{#1}}}
% Dato que el grupo debe completar antes de entregar. Buscar "PENDIENTE".
\newcommand{\pendiente}[1]{\textcolor{cPend}{\textbf{[PENDIENTE: #1]}}}
\newcolumntype{L}[1]{>{\raggedright\arraybackslash}p{#1}}
\newcolumntype{Y}{>{\raggedright\arraybackslash}X}

% ---------------------------------------------------------------------------
\begin{document}

\begin{titlepage}
\centering
{\large\bfseries Instituto Tecnológico de Costa Rica\par}
\vspace{0.3cm}
{\large Escuela de Ingeniería en Computadores\par}
\vspace{0.3cm}
{\large CE-1113 Sistemas Empotrados\par}
\vspace{3.2cm}
{\LARGE\bfseries Documento de Aprendizaje Continuo (AC)\par}
\vspace{0.6cm}
{\Large Proyecto I --- Sistema embebido a la medida para un robot aspiradora
autónomo con reproducción de audio y control remoto\par}
\vspace{3.2cm}
{\large\bfseries Integrantes\par}
\vspace{0.2cm}
{\large \pendiente{nombres y carnés de los cuatro integrantes}\par}
\vspace{1.6cm}
{\large\bfseries Profesor\par}
\vspace{0.2cm}
{\large Dr.-Ing. Jeferson González Gómez\par}
\vfill
{\large Octubre de 2026\par}
\end{titlepage}

\tableofcontents
\newpage

% ===========================================================================
\section{Introducción}

Este documento evidencia el atributo de \textbf{Aprendizaje Continuo (AC)} en
el Proyecto I. Identifica lo que el grupo no sabía y necesitaba saber (AC1),
las tecnologías que contribuyeron a aprenderlo (AC2), las estrategias
concretas que se aplicaron (AC3) y una evaluación de qué tan bien funcionaron
(AC4).

Cada afirmación se apoya en un rastro verificable del repositorio: un
documento de \id{docs/}, un commit o un issue. El aprendizaje que aquí se
describe es el que dejó evidencia.

% ===========================================================================
\section{AC1 --- Necesidades de aprendizaje}

\subsection{Contexto}

El proyecto pide construir un sistema operativo a la medida, una biblioteca
de control, un servidor web y un modelo físico con su electrónica de
potencia. Al inicio, el grupo tenía experiencia en programación en C, en
Linux como usuario y en electrónica digital básica. No había construido una
imagen con Yocto, ni diseñado una etapa de potencia aislada, ni integrado un
sensor inercial.

A eso se suma el cambio tecnológico. Las herramientas del proyecto cambian
más rápido que la documentación que las explica: la sintaxis de Yocto varió
entre versiones, las distribuciones del host son más nuevas que las validadas
por la rama \id{scarthgap}, y el kernel 6.6 exige un módulo de WiFi que
versiones anteriores no pedían. Buena parte de lo que se encuentra en línea
está escrito para una versión que ya no aplica, y saber distinguirlo resultó
ser una necesidad de aprendizaje en sí misma.

\subsection{Conocimientos}

\begin{longtable}{L{3.4cm}L{5.6cm}L{6.0cm}}
\toprule
\textbf{Área} & \textbf{Qué hacía falta aprender} & \textbf{Por qué lo exigía el proyecto} \\
\midrule
\endhead
Yocto y BitBake & Capas, recetas, \id{bbappend}, variables con
\emph{overrides}, caché \id{sstate}, tablas de particiones con \id{wic}. &
La imagen mínima y la receta propia son obligatorias. \\
Compilación cruzada & Toolchain-SDK, \emph{sysroot}, archivo de toolchain
de CMake, verificación con \id{file} y \id{readelf}. & Nada puede compilarse
en el target. \\
\id{systemd} & Unidades, dependencias (\id{Requires} frente a \id{Wants}),
política de reinicio, objetivos de red. & Arranque automático con
\id{Restart=on-failure}. \\
Audio en Linux embebido & ALSA, \emph{device tree overlays}, la salida PWM
de la Raspberry Pi. & Reproducir MP3 por un parlante a bordo. \\
Temporización de GPIO & PWM y pulsos de servo por DMA, medición de pulsos
en microsegundos, I2C. & Motores, servo, HC-SR04 y MPU-6050. \\
Electrónica de potencia & Aislamiento óptico, separación de tierras,
tiempos de conmutación de un optoacoplador, caída de un puente bipolar. & La
nota de seguridad del enunciado es de lectura obligatoria. \\
Baterías de Li-Ion & Rango de tensión de una celda, función del BMS,
dimensionamiento de un reductor. & Fuente portátil sin dañar el equipo
prestado. \\
Estimación de estado & Integración de un acelerómetro y su deriva, filtro
complementario, corrección de velocidad cero. & Odometría y mapa sin
encoders. \\
Servidor HTTP en C & \id{libmicrohttpd}, sesiones, resumen de contraseñas. &
Control remoto con autenticación. \\
\bottomrule
\end{longtable}

\subsection{Habilidades y destrezas}

\begin{itemize}
  \item \textbf{Diagnosticar sin pantalla.} El robot no tiene monitor ni
        teclado: todo se depura por la consola serie, por SSH y con
        \id{journalctl}.
  \item \textbf{Leer registros largos de construcción} y encontrar la línea
        que importa entre miles.
  \item \textbf{Medir antes de conectar.} Usar el multímetro como paso
        obligatorio, no como último recurso.
  \item \textbf{Soldar y montar} una etapa de potencia en placa perforada.
  \item \textbf{Trabajar con un flujo de Git profesional}: ramas, Pull
        Requests, Conventional Commits, issues e hitos.
\end{itemize}

\subsection{Actitudes}

\begin{itemize}
  \item \textbf{Desconfiar del mensaje de error.} Varias veces el mensaje
        apuntaba a una causa falsa.
  \item \textbf{Paciencia con ciclos lentos.} La primera construcción de la
        imagen toma unas tres horas; un error de configuración se paga caro.
  \item \textbf{Responsabilidad sobre equipo ajeno.} Un cable mal puesto
        destruye una tarjeta que no es del grupo.
  \item \textbf{Documentar mientras se trabaja}, no al final.
\end{itemize}

% ===========================================================================
\section{AC2 --- Tecnologías nuevas y emergentes}

Estas son las tecnologías que el grupo no había usado antes y que
contribuyeron al aprendizaje. Para cada una se indica qué permitió aprender.

\begin{longtable}{L{3.6cm}L{11.4cm}}
\toprule
\textbf{Tecnología} & \textbf{Cómo contribuyó al aprendizaje} \\
\midrule
\endhead
Yocto Project, versión 5.0 (\id{scarthgap}) & Obligó a entender de qué está
hecho un sistema Linux: cada paquete de la imagen es una decisión explícita.
Justificar cada uno (\id{docs/paquetes.md}) fue la mejor forma de aprender
qué hace cada uno. \\
Emulación con QEMU & En modo usuario (\id{qemu-aarch64}) permitió ejecutar
los binarios ARM en la laptop. En modo sistema, con una máquina propia
(\id{qemuarm64-robot}), permitió arrancar la imagen completa y observar
\id{systemd} sin la Raspberry Pi. \\
Simulación del hardware & Un reemplazo de \id{pigpiod} que mueve un robot
virtual (\id{sim/}). Escribirlo exigió entender con precisión qué hace cada
llamada de GPIO, y permitió experimentar con la navegación sin riesgo. \\
Información de presión del kernel (PSI) & \id{/proc/pressure} y las variables
\id{BB_PRESSURE_MAX_*} de BitBake: una forma reciente de regular la carga por
la presión real de memoria, en lugar de un número fijo de hilos. \\
Sensores inerciales MEMS & El MPU-6050 llevó a estudiar fusión de sensores
con un caso real y medible. \\
\id{pigpio} con DMA & Mostró cómo generar temporización estable desde espacio
de usuario en un sistema que no es de tiempo real. \\
Amplificación clase D & El PAM8403 frente al LM386: el mismo problema
resuelto con una fracción del consumo. \\
\emph{Device tree overlays} & \id{audremap} mostró que redirigir un
periférico a otro pin es configuración, no código. \\
GitHub Issues, hitos, tablero de proyecto y la CLI \id{gh} & Hicieron visible
el estado del proyecto para los cuatro integrantes y dejaron el rastro de
cada decisión. \\
LaTeX con TikZ y CircuiTikZ & Los diagramas de hardware viven como texto en
el repositorio y se revisan como cualquier otro cambio. \\
\bottomrule
\end{longtable}

\pendiente{si el grupo usó asistentes de inteligencia artificial generativa
u otras tecnologías emergentes, describir aquí para qué se usaron y cómo se
verificó lo que produjeron}

% ===========================================================================
\section{AC3 --- Estrategias implementadas}

\subsection{Ir a la fuente primaria}

Ante un problema, la primera consulta fue la documentación oficial, la hoja
de datos o el código fuente, antes que un foro.

El caso más claro fue el fallo de descarga del kernel. BitBake reportaba que
una revisión no existía en la rama. Leyendo el código del \emph{fetcher} se
encontró que ejecuta \id{git branch --contains} descartando la salida de
error, y que las versiones recientes de git abortan ese comando en el
repositorio espejo. El mensaje era falso; la causa y la solución quedaron
documentadas en \id{docs/yocto-setup.md}.

\subsection{Construir de menor a mayor}

Cada capa se validó antes de montar la siguiente: primero
\id{core-image-minimal}, para comprobar el BSP, el WiFi y el audio; luego la
imagen del proyecto. En el hardware se siguió el mismo orden, con los
procedimientos de cinco pasos de \id{docs/hardware-aislamiento.md} y
\id{docs/hardware-alimentacion.md}: aislamiento sin energía, riel de potencia
solo, optoacopladores con fuente de banco, riel lógico solo, y la Raspberry
Pi de última.

\subsection{Simular para no depender del hardware}

El hardware llegó después que el software. Para no detener el aprendizaje se
construyeron tres ambientes sin la Raspberry Pi: la biblioteca contra el
simulador en la laptop, el binario ARM bajo \id{qemu-aarch64}, y la imagen
completa en QEMU. El programa \id{sim/prueba_librobot} ejercita cada módulo
con 47 comprobaciones.

\subsection{Escribir para entender}

Cada decisión de diseño se documentó en el momento de tomarla, con la
alternativa descartada y el motivo. La carpeta \id{docs/} tiene trece
documentos técnicos escritos durante el desarrollo, no después. Explicar por
escrito por qué 390\,$\Omega$ y no 220\,$\Omega$ obliga a hacer la cuenta.

\subsection{Organizar el tiempo}

\begin{itemize}
  \item El enunciado se desglosó en 37 issues agrupados en 10 hitos con
        fecha, de \emph{Repositorio y flujo de trabajo} (8 de setiembre) a
        \emph{Documentación y entrega final} (6 de octubre).
  \item Un tablero de proyecto con tres columnas mostró en todo momento qué
        estaba pendiente, en curso y terminado.
  \item Cada integrante trabajó en su rama; los cambios entraron a
        \id{develop} y a \id{main} por Pull Request (20 integrados).
\end{itemize}

\subsection{Medir antes y después}

Los problemas de desempeño se trataron como experimentos. El servidor
arrancaba a los 130\,s; se identificaron dos causas independientes en la
configuración de red, se corrigieron y se volvió a medir: 17\,s. El commit
registra la medición anterior, la causa y la medición posterior.

\subsection{Revisar entre pares y con herramientas}

Los Pull Requests fueron el punto de revisión entre integrantes. Se sumaron
verificaciones automáticas donde existían: \id{systemd-analyze verify}
detectó una directiva ubicada en la sección equivocada que \id{systemd}
ignoraba sin avisar.

\subsection{Búsqueda bibliográfica}

Las hojas de datos del PC817, el L298N, el MPU-6050, el PAM8403 y el MP2307
fueron la base de los cálculos de \id{docs/hardware-*.md}: corriente del LED
del optoacoplador, relación de transferencia de corriente, caída del puente,
ganancia del amplificador.

% ===========================================================================
\section{AC4 --- Evaluación crítica de las estrategias}

\begin{longtable}{L{2.9cm}L{5.9cm}L{6.2cm}}
\toprule
\textbf{Estrategia} & \textbf{Qué funcionó} & \textbf{Qué no funcionó} \\
\midrule
\endhead
Ir a la fuente primaria & Resolvió problemas que ninguna búsqueda habría
resuelto, como el del \emph{fetcher}. Lo aprendido se sostiene al cambiar de
versión. & Es lenta. Leer el código de BitBake tomó horas que un foro, de
haber tenido la respuesta, habría ahorrado. \\
Construir de menor a mayor & Ningún componente se dañó. Cada fallo quedó
acotado a la última capa agregada. & Pospone los problemas de integración:
los que solo aparecen con todo junto se descubren tarde. \\
Simular & El software avanzó semanas antes de tener el robot: biblioteca,
servidor, navegación y mapa se probaron sin hardware. & El simulador solo
contiene lo que el grupo supo modelar. No modela el tiempo de apagado del
optoacoplador, y por eso la PWM de 1\,kHz funcionó en la simulación y no
movió los motores reales. \\
Escribir para entender & Detectó errores antes de armar: la inversión del
optoacoplador se encontró al documentar el circuito, no al encenderlo. & Cada
cambio de diseño obligó a actualizar varios documentos. Cuando el radar
reemplazó a los tres sensores, parte de la documentación quedó atrasada un
tiempo. \\
Organizar el tiempo & El tablero evitó trabajo duplicado y dejó claro qué
faltaba. & Las fechas de los hitos fueron optimistas: cinco vencieron con
issues abiertos. No se previó cuánto dependía cada hito de tener el hardware
armado. \\
Medir antes y después & Convirtió una impresión en un dato: 130\,s a 17\,s.
& Las métricas de eficiencia se dejaron para el final, cuando ya había poco
margen para actuar sobre ellas. \\
Revisión entre pares & Los Pull Requests dejaron el rastro de cada cambio. &
Bajo presión de tiempo, varios se aprobaron sin una revisión a fondo. \\
\bottomrule
\end{longtable}

\subsection{Lo que el grupo haría distinto}

\begin{enumerate}
  \item \textbf{Encender el hardware antes.} Una prueba mínima de un motor a
        través de un optoacoplador en la primera semana habría mostrado el
        problema de la PWM con un mes de margen. La simulación complementa a
        la placa, no la reemplaza.
  \item \textbf{Aclarar las ambigüedades del enunciado al inicio.} La
        decisión de cubrir las detecciones frontal y lateral con un solo
        sensor sobre un servo dependía de una interpretación que convenía
        confirmar con el profesor antes de diseñar alrededor de ella.
  \item \textbf{Preparar el host de construcción desde el primer día.} Los
        cuelgues por falta de memoria costaron días que una máquina adecuada,
        o la regulación por presión desde el inicio, habrían evitado.
  \item \textbf{Medir de forma continua.} Tomar las métricas con cada imagen
        permite reaccionar; tomarlas al final solo permite reportarlas.
  \item \textbf{Poner fechas a los hitos según sus dependencias}, no según
        el calendario.
\end{enumerate}

\subsection{Balance}

Las estrategias cumplieron su propósito: el grupo terminó el proyecto
dominando herramientas que al inicio no conocía, y lo aprendido quedó escrito
de forma que otro grupo podría reproducirlo. La lección que se repite es que
las estrategias que más enseñaron ---leer la fuente, documentar, simular---
son también las más costosas en tiempo, y que ninguna sustituye el contacto
temprano con el hardware real.

\pendiente{revisar el documento entre los cuatro integrantes, en especial
el punto de partida de la sección 2.1 y esta evaluación: deben reflejar la
experiencia real del grupo}

% ===========================================================================
\section{Conclusiones}

\begin{itemize}
  \item Las necesidades de aprendizaje más grandes no estaban en la
        programación sino en lo que la rodea: construir el sistema
        operativo, aislar la electrónica de potencia y diagnosticar sin
        pantalla.
  \item La emulación y la simulación fueron las tecnologías que más
        contribuyeron, porque permitieron aprender sin esperar al hardware y
        sin arriesgarlo.
  \item Aprender en un entorno que cambia exige ir a la fuente primaria y
        verificar la versión de todo lo que se lee.
  \item La evaluación crítica muestra un límite claro: lo que no se modela
        no se aprende hasta tocar el hardware.
\end{itemize}

\end{document}
